\documentclass[11pt]{article}
\usepackage{amsmath,amssymb,geometry}
\geometry{margin=22mm}
\usepackage[UTF8]{ctex}
\begin{document}

\section*{P353：确定性多模式铁路货运模型}
这是当前为某个供应需求情景所执行的完整模型M0，其中包含了明确的OD和截止日期相关的修复措施。这是一个混合整数二次约束问题：产品$\alpha\eta$和$\delta\theta$被精确保留，并通过全局非凸优化方法进行处理。每个需求情景都有各自的决策，并且可以独立求解。

\paragraph{集合与固定参数。}
$Y$是OD对集合，$R$是货运列车集合，$K$是火车集合。  
$K_y^1,K_y^2,K_y^3$是适用于$y$的兼容式样列车、预留车厢以及专用货运列车。每列当前运行的列车都属于其指定的OD。  
$O,D,N$是独立的起点角色、终点角色以及中间站集合。  
$I=O\cup D\cup N$中，同一物理站点的不同角色之间有明确的距离/时间映射关系。  
所有$p,q,a,u^1,u^2,u_3,u_4,w,l,c_1,c_2,c_3,c_4,c_5,c_6,
d_1,d_2,t_1,t_2,t_3,t_4,t_5,t_6,t_7,M$都是数值参数，其索引与以下公式中的索引完全一致。  
特别是$c_1,c_6,t_5$是标量，$a_{yi}$是OD端点的关联信息，$l_{kr}$是固定列车预定到达时间的标志，$t_1^r$是货运的可用性，而$t_7^r$则是其截止日期。  
源代码中提到的正数值域是特意设定的：$\beta_k\in\{1,2,3\}$和$\iota\in\mathbb Z_{\ge1}$，其中正数值域也是明确指定的。  
专用列车路径变量遵循所指定的规则，适用于使用过的路径和未使用的路径。
\subsection*{决策变量}
\[
\begin{aligned}
& \alpha_{ykr} && \text{quantity of freights of type $r$ transported by high-speed train $k$ of OD pair $y$.} \\
& \beta_{k} && \text{quantity of reserved coaches for reserved-coach train $k$, $k \in K^2_y$.} \\
& \gamma_{k} && \text{binary; equals one if freight train $k$ is arranged for express delivery, otherwise zero, $k \in K^3_y$.} \\
& \delta_{ykr} && \text{binary; equals one if freights of type $r$ are transported by high-speed train $k$ of OD pair $y$, otherwise zero.} \\
& \zeta_{ykri} && \text{binary; equals one if OD pair $y$'s freight train $k$, which transports freights of type $r$, passes through station $i$, otherwise zero, $k \in K^3_y$.} \\
& \eta_{ykrii'} && \text{binary; equals one if OD pair $y$'s freight train $k$, which transports freights of type $r$, passes through station $i'$ after passing through station $i$, otherwise zero, $k \in K^3_y$.} \\
& \theta_{k} && \text{departure time of freight train $k$, $k \in K^3_y$.} \\
& \iota_{ykri} && \text{sequence number of station $i$, through which OD pair $y$'s freight train $k$ passes, $k \in K^3_y$; the train transports freights of type $r$.}
\end{aligned}
\]

\subsection*{目标函数}
\[
\begin{aligned}
\text{Maximize} \quad & \sum_{y \in Y} \sum_{k \in K} \sum_{r \in R} p_r \alpha_{ykr} \\
& - \sum_{y \in Y} \sum_{k \in K} \sum_{r \in R} 2c_1 \alpha_{ykr} \\
& - \sum_{y \in Y} \sum_{k \in K^2_y} \sum_{r \in R} c_2^k d_1^k \alpha_{ykr} \\
& - \sum_{y \in Y} \sum_{k \in K^2_y} c_3 \beta_k \\
& - \sum_{y \in Y} \sum_{k \in K^3_y} \sum_{r \in R} \sum_{i \in N \cup O} \sum_{i' \in N \cup D} c_2^k d_2^{ii'} \alpha_{ykr} \eta_{ykrii'} \\
& - \sum_{y \in Y} \sum_{k \in K^3_y} c_4^k \gamma_k \\
& - \sum_{y \in Y} \sum_{r \in R} c_5^r \left( q_{yr} - \sum_{k \in K} \alpha_{ykr} \right) \\
& - \sum_{y \in Y} \sum_{k \in K^1_y \cup K^2_y} \sum_{r \in R} c_6 \delta_{ykr} (t_2^k - t_1^r) \\
& - \sum_{y \in Y} \sum_{k \in K^3_y} \sum_{r \in R} c_6 \delta_{ykr} (\theta_k - t_1^r)
\end{aligned}
\]

\subsection*{约束条件：}
\[
\begin{aligned}
& \sum_{k \in K} \alpha_{ykr} \leq q_{yr} && \forall y \in Y, r \in R && (2) \\
& \alpha_{ykr} \leq M \delta_{ykr} && \forall y \in Y, k \in K, r \in R && (3) \\
& \sum_{y \in Y} \sum_{k \in K} \sum_{r \in R} a_{yi} \alpha_{ykr} \leq u_i^1 && \forall i \in O \cup D && (4) \\
& \sum_{r \in R} \alpha_{ykr} \leq u_k^2 && \forall y \in Y, k \in K^1_y \cup K^3_y && (5) \\
& \sum_{r \in R} a_{yi} \alpha_{ykr} \leq t_{ki}^3 w_i && \forall y \in Y, k \in K^1_y \cup K^2_y, i \in O \cup D && (6) \\
& \delta_{ykr} (t_1^r - t_2^k) \leq 0 && \forall y \in Y, k \in K^1_y \cup K^2_y, r \in R && (7) \\
& (t_2^k + t_4^k + 2 t_5 \alpha_{ykr}) l_{kr} \leq t_7^r && \forall y \in Y, k \in K^1_y \cup K^2_y, r \in R && (8) \\
& \delta_{ykr} \leq \beta_k && \forall y \in Y, k \in K^2_y, r \in R && (9) \\
& \beta_k \leq 3 && \forall y \in Y, k \in K^2_y && (10) \\
& \sum_{r \in R} \alpha_{ykr} \leq u_3 \beta_k && \forall y \in Y, k \in K^2_y && (11) \\
& \delta_{ykr} \leq \gamma_k && \forall y \in Y, k \in K^3_y, r \in R && (12) \\
& \delta_{ykr} (t_1^r - \theta_k) \leq 0 && \forall y \in Y, k \in K^3_y, r \in R && (13) \\
& \sum_{i' \in N} \eta_{ykrii'} = a_{yi} && \forall y \in Y, k \in K^3_y, r \in R, i \in O && (14) \\
& \sum_{i \in N} \eta_{ykrii'} = a_{yi'} && \forall y \in Y, k \in K^3_y, r \in R, i' \in D && (15) \\
& \sum_{i' \in N \cup O} \eta_{ykri'i} = \zeta_{ykri} && \forall y \in Y, k \in K^3_y, r \in R, i \in N && (16) \\
& \sum_{i' \in N \cup D} \eta_{ykrii'} = \zeta_{ykri} && \forall y \in Y, k \in K^3_y, r \in R, i \in N && (17) \\
& \iota_{ykri} - M (1 - \eta_{ykrii'}) + 1 \leq \iota_{ykri'} && \forall y \in Y, k \in K^3_y, r \in R, i, i' \in N && (18) \\
& \theta_k + \sum_{i \in N \cup O} \sum_{i' \in N \cup D} t_6^{ii'} \eta_{ykrii'} + 2 t_5 \alpha_{ykr} \leq t_7^r && \forall y \in Y, k \in K^3_y, r \in R && (19) \\
& \sum_{y \in Y} \sum_{k \in K^3_y} \sum_{r \in R} \alpha_{ykr} \eta_{ykrii'} \leq u_4^{ii'} && \forall i, i' \in I && (20) \\
& \alpha_{ykr} \geq 0 && \forall y \in Y, k \in K, r \in R && (21) \\
& \beta_k \in \mathbb{N}^+ && \forall y \in Y, k \in K^2_y && (22) \\
& \delta_{ykr} \in \{0, 1\} && \forall y \in Y, k \in K, r \in R && (23) \\
& \gamma_k, \zeta_{ykri}, \eta_{ykrii'} \in \{0, 1\} && \forall y \in Y, k \in K^3_y, r \in R, i, i' \in I && (24) \\
& \theta_k \geq 0 && \forall y \in Y, k \in K^3_y && (25) \\
& \iota_{ykri} \in \mathbb{N}^+ && \forall y \in Y, k \in K^3_y, r \in R, i \in I && (26)
\end{aligned}
\]


\subsection*{明确的资格修复措施}
除了上述每个显示的目标术语、约束条件和领域之外，还需施加…
\[
\alpha_{ykr}=\delta_{ykr}=0
\quad(k\notin K_y^1\cup K_y^2\cup K_y^3),
\]
\[
\alpha_{ykr}=\delta_{ykr}=0
\quad(k\in K_y^1\cup K_y^2,\ l_{kr}=0).
\]
在已完成的输入数据中，
$l_{kr}=0$与$t_2^k+t_4^k>t_7^r$的时间点完全一致。
这些是对提取出的模型中的遗漏部分的修复，得到了源代码的OD列车定义和截止日期要求的支持（第7-8页）；
修复后的索引定义将在整个模型中继续使用。
所有原始行（2–26）都被保留下来，包括行（8）中修复后的$l_{kr}$。

\paragraph{实施与证据。}
当前代码中的第23行到第49行实现了输入轴、域以及所有目标项；第50行到第76行实现了(2)--(26)部分的功能；而第77行到第84行则实现了两次修复操作。这些代码会独立地评估所有线性约束和二次约束、所有变量域以及目标值。每个记录的场景都有其自己的目标和结果；而加权汇总则将这些特定于场景的结果整合在一起。

\end{document}
